\documentclass[12pt]{letter}

\usepackage{geometry}
\geometry{margin=1in}
\usepackage{hyperref}

\begin{document}

\begin{letter}{The Editors\\
\textit{Foundations of Physics}\\
Springer Nature}

\opening{Dear Editors,}

I submit for your consideration the manuscript entitled:

\begin{center}
\textbf{Constraint-Induced Interface Realism: A Complete Axiomatic Framework\\
for Bounded Cognitive Systems and Observable Reality Reconstruction}
\end{center}

\textbf{Problem addressed.}
No unified formal framework has previously established, with the axiomatic
completeness and mathematical rigor comparable to Kolmogorov probability theory
or Hilbert-space quantum mechanics, the precise relationship between structural
constraints on bounded cognitive systems and the interface representations such
systems construct as operational models of reality.  Existing approaches either
provide informal philosophical accounts or import quantum formalism without
rigorously deriving which axioms are required and which are redundant.

\textbf{Contributions of this manuscript.}
This paper presents the theory of Constraint-Induced Interface Realism (CIIR)
as a self-contained formal system $\mathcal{T}_\mathrm{CIIR} = (\mathcal{S},
\mathcal{A}, \mathcal{R}, \mathcal{D}, \mathcal{M}, \mathcal{I})$ consisting
of a state space, axiom system, operator algebra, derivation system, model
class, and interpretation map.  The specific contributions are:

\begin{enumerate}
\item A minimal, mutually independent set of six axioms whose independence is
      formally demonstrated by explicit counter-models for each.
\item A derivation of the full operator algebra, including constraint
      commutation relations tied to the curvature of a Riemannian constraint
      manifold.
\item The CIIR Representation Theorem establishing a functorial Hilbert-space
      representation for every bounded cognitive system.
\item A canonical Gorini--Kossakowski--Sudarshan--Lindblad master equation
      governing constrained cognitive evolution, with stability analysis.
\item Ten non-trivial theorems with formal proofs, including the
      Constraint-Induced Distortion Theorem, Interface Non-Invertibility
      Theorem, Cognitive Interference Theorem, and a CIIR N\"{o}ther Theorem.
\item Explicit construction of finite-dimensional (qubit) and
      infinite-dimensional (harmonic oscillator) model classes, each verified
      to satisfy all axioms.
\item Three quantitatively falsifiable predictions expressed as testable
      inequalities relating curvature, cognitive volume, and observable
      quantities (decision latency, decoherence rates, channel capacity).
\item Embedding of classical Kolmogorov probability, Bayesian inference, and
      quantum cognition (Busemeyer--Bruza) as limiting cases.
\end{enumerate}

\textbf{Why this manuscript deserves review.}
The framework is presented at the level of rigor expected of graduate-level
mathematical physics.  Every symbol is defined before use; all theorems are
stated precisely and proved (or rigorously sketched); all axioms are
independently justified.  The theory is non-trivial (admits rich models),
internally consistent (no contradictions derive from the axioms), and
operationalizable (predictions are expressed as measurable quantities).
CIIR addresses foundational questions in the philosophy of science, cognitive
science, and quantum information theory, making it suitable for the
\textit{Foundations of Physics} readership.

\textbf{Conflicts of interest.}
None.

\textbf{Data availability.}
This is a theoretical manuscript; no experimental data are reported.  All
mathematical derivations are contained within the manuscript.

\closing{Sincerely,}

Robert Ringler\\
Independent Researcher\\
\href{https://github.com/robertringler/QRATUM}{https://github.com/robertringler/QRATUM}

\end{letter}
\end{document}
